% ===== PRÉAMBULE CANONIQUE — à coller VERBATIM en tête de CHAQUE .tex =====
\documentclass[11pt,a4paper]{article}
\usepackage[utf8]{inputenc}\usepackage[T1]{fontenc}\usepackage{lmodern}
\usepackage[french]{babel}
\usepackage{amsmath,amssymb,mathtools}\usepackage{icomma}
\usepackage{siunitx}\sisetup{locale=FR}  % \qty{}{}, \si{}, \ang{}
\usepackage{xcolor}\usepackage{colortbl}
\usepackage{tikz}\usepackage{pgfplots}\pgfplotsset{compat=1.18}
\usepackage[siunitx,europeanresistors]{circuitikz} % circuits (élec.)
\usepackage[most]{tcolorbox}\usepackage{enumitem}\usepackage{array,booktabs}
\usepackage[a4paper,margin=2.2cm]{geometry}
\usetikzlibrary{arrows.meta,calc,angles,quotes,positioning,patterns,%
  backgrounds,shapes.geometric,decorations.pathmorphing,decorations.markings,intersections}
\definecolor{primary}{RGB}{222,49,99}\definecolor{primarydark}{RGB}{150,22,55}
\definecolor{defcol}{RGB}{0,114,178}\definecolor{methcol}{RGB}{52,92,148}
\definecolor{excol}{RGB}{0,135,90}\definecolor{attncol}{RGB}{200,40,55}
\tcbset{boxbase/.style={enhanced,breakable,boxrule=0.9pt,arc=2.5pt,
  left=3mm,right=3mm,top=2.3mm,bottom=2.3mm,fonttitle=\bfseries,coltitle=white}}
% --- boîtes TUTORIEL (noms EXACTS, ne pas renommer) ---
\newtcolorbox{cle}[1][]{boxbase,colback=defcol!6,colframe=defcol,
  title={\textsc{À retenir}\ifx\relax#1\relax\relax\else\textmd{ : \itshape #1}\fi}}
\newtcolorbox{methode}[1][]{boxbase,colback=methcol!7,colframe=methcol,sharp corners,
  title={\textsc{Méthode}\ifx\relax#1\relax\relax\else\textmd{ --- \itshape #1}\fi}}
\newtcolorbox{exemplebox}[1][]{boxbase,colback=excol!5,colframe=excol,
  title={\textsc{Exemple}\ifx\relax#1\relax\relax\else\textmd{ : \itshape #1}\fi}}
\newtcolorbox{piege}[1][]{boxbase,colback=attncol!6,colframe=attncol,
  title={\textsc{Piège}\ifx\relax#1\relax\relax\else\textmd{ : \itshape #1}\fi}}
\newtcolorbox{remarque}[1][]{boxbase,colback=black!4,colframe=black!45,
  title={\textsc{Remarque}\ifx\relax#1\relax\relax\else\textmd{ : \itshape #1}\fi}}
% --- boîtes EXERCICE / CORRIGÉ (noms EXACTS, auto-comptées) ---
\newtcolorbox[auto counter]{exo}[1][]{boxbase,colback=primary!4,colframe=primary,
  title={\textsc{Exercice}~\thetcbcounter\ifx\relax#1\relax\relax\else\textmd{ --- \itshape #1}\fi}}
\newtcolorbox[auto counter]{corrige}[1][]{boxbase,colback=excol!4,colframe=excol,
  title={\textsc{Corrigé}~\thetcbcounter\ifx\relax#1\relax\relax\else\textmd{ --- \itshape #1}\fi}}
\newcommand{\vv}[1]{\vec{#1}}\newcommand{\ds}{\displaystyle}
\newcommand{\R}{\mathbb{R}}\newcommand{\N}{\mathbb{N}}\newcommand{\Z}{\mathbb{Z}}
% =========================================================================
\begin{document}
\sisetup{per-mode=symbol}
% ================= FACILE =================

\begin{exo}[énergie cinétique]
Un satellite de masse $m=\qty{500}{\kilogram}$ se déplace à la vitesse
$v=\qty{7000}{\metre\per\second}$. Calcule son énergie cinétique $E_c=\tfrac12 m v^2$.
\end{exo}

\begin{exo}[énergie potentielle]
Le même satellite ($m=\qty{500}{\kilogram}$) est à la distance $r=\qty{1.0e7}{\metre}$
du centre de la Terre ($\mathcal{G}M_T=\qty{4.0e14}{}$ SI). Calcule son énergie
potentielle $E_p=-\dfrac{\mathcal{G}M_T m}{r}$ (attention au signe).
\end{exo}

\begin{exo}[énergie mécanique]
Sur une orbite circulaire, l'énergie mécanique vérifie $E_m=-E_c$. Un satellite a une
énergie cinétique $E_c=\qty{3.0e10}{\joule}$. Que vaut son énergie mécanique $E_m$ ?
\end{exo}

\begin{exo}[le signe des énergies]
Pour un satellite lié à la Terre, indique le signe (positif ou négatif) de :
\begin{enumerate}[label=\alph*)]
  \item l'énergie cinétique $E_c$ ;
  \item l'énergie potentielle $E_p$ ;
  \item l'énergie mécanique $E_m$.
\end{enumerate}
\end{exo}

\begin{exo}[tout à partir de $E_c$]
Sur une orbite circulaire, $E_p=-2E_c$ et $E_m=-E_c$. Un satellite a
$E_c=\qty{1.5e10}{\joule}$.
\begin{enumerate}[label=\alph*)]
  \item Donne son énergie potentielle $E_p$.
  \item Donne son énergie mécanique $E_m$.
\end{enumerate}
\end{exo}

\end{document}
